\begin{corollary}[The unconditional K3 range]
\label{c:k3-stable-range}
\status{Proved}
For every projective complex K3 surface with its Beauville--Voisin class, the requested relations hold in every power $S^k$ with $k\leq45$. More generally they hold in normalized degree $d$ of every $S^k$ whenever
\[
 \min(d,k-d)\leq22,
\]
where normalized degree $d$ means Chow codimension $2d$. In this range no finite-dimensionality assumption is required.
\end{corollary}

\begin{proof}
The Beauville--Voisin identities give the local relations, so every Chow monomial reduces to a linear combination of disjoint matching monomials
$\prod_{i\in I}o_i\prod_{\{a,b\}\in\mathcal M}q_{ab}$.
If the matching has $m$ edges, degree $d$ means $|I|+m=d$, and the number of used factors is $|I|+2m$. Consequently $m\leq\min(d,k-d)$. For $m\leq22$, the orthogonal matching tensors in $(H^2(S))^{\otimes2m}$ are linearly independent: the first determinantal relation requires $23$ edges, as proved in Theorem~\ref{c:pair-ranks}. Monomials with different point and matching supports lie in different K\"unneth multidegrees. Thus their cohomological classes are independent, and no homologically zero Chow combination remains. For $k\leq45$, every matching has at most $22$ edges. This proves both assertions.
\end{proof}

The remaining relation cannot be obtained merely by asserting that all lower-power relations hold. It first appears in the middle normalized degree of $S^{46}$ and is exactly \eqref{c:k3-determinant}. We also tested three possible ways to remove this obstruction. A Kuga--Satake correspondence giving a splitting in homological motives would produce a composite $1+\eta$ in Chow, with $\eta$ homologically trivial; promoting it to a Chow splitting still requires nilpotence or another argument making $1+\eta$ invertible. Specializing to dense loci with finite-dimensional motives does not prove the generic Chow vanishing, because the rational-equivalence locus is only controlled as a countable union of closed loci without a uniform family of equivalences. Finally, the known small-diagonal identities are precisely the local matching reductions and do not supply the first orthogonal determinant. These are failed proof steps, not counterexamples to finite-dimensionality.
